% Part: model-theory
% Chapter: lindstrom
% Section: lindstrom-proof

\documentclass[../../../include/open-logic-section]{subfiles}

\begin{document}

\olfileid{mod}{lin}{prf}

\olsection{مبرهنة ليندستروم}

\begin{lem}
\ollabel{lem:lindstrom}
لتكن $!E \in {\OLId{latin-looped}{L}}(\Lang{L})$، ولتكن $\Lang{L}$ منتهية. فإن
وجد ${\OLId{latin-ordinary}{n}} \in \Nat$ تتحقق معه، لكل !!{structure}s اثنتين $\Struct{M}$ و
$\Struct{N}$، القضية الآتية: إذا كان $\Struct{M} \equiv_{\OLId{latin-ordinary}{n}} \Struct{N}$ و
$\Struct{M} \models_{\OLId{latin-looped}{L}} !E$، فإن $\Struct{N} \models_{\OLId{latin-looped}{L}} !E$ أيضًا؛
فإن $!E$ تكافئ~!!{sentence} من الرتبة الأولى؛ أي توجد $!D$ من الرتبة
الأولى بحيث $\Mod({\OLId{latin-looped}{L}}){!E}=\Mod({\OLId{latin-looped}{L}}){!D}$.
\end{lem}

\begin{proof}
ليكن ${\OLId{latin-ordinary}{n}}$ بحيث تتفق أي !!{structure}s متكافئتين حتى~${\OLId{latin-ordinary}{n}}$، $\Struct{M}$
و$\Struct{N}$، على القيمة المسندة إلى~$!E$. وقد ثبت في
\olref[bas][pis]{prop:qr-finite} أن !!{language} منتهية لا تشتمل، إذا سوّينا بين المتكافئات
منطقيًا، إلا على عدد منتهٍ من !!{sentence}s الرتبة الأولى التي لا
تتجاوز رتبة كمها~${\OLId{latin-ordinary}{n}}$. ولكل !!{structure} ثابتة~$\Struct{M}$، لتكن
$!D_{\Struct{M}}$ اقتران ممثل واحد من كل فئة تكافؤ منطقي للـ!!{sentence}s
$!A$ من الرتبة الأولى الصادقة في~$\Struct{M}$ والتي تحقق
$\QuantRank{!A} \le {\OLId{latin-ordinary}{n}}$ (فهذا الاقتران منتهٍ)، بحيث
$\Struct{N} \models !D_{\Struct{M}}$ إذا وفقط إذا كان
$\Struct{N} \equiv_{\OLId{latin-ordinary}{n}} \Struct{M}$. ثم ضع
$!D = \textstyle\bigvee
\Setabs{!D_{\Struct{M}}}{\Struct{M} \models_{\OLId{latin-looped}{L}} !E}$؛ وهذا الانفصال منتهٍ
أيضًا حتى التكافؤ المنطقي.

فنحصل على $\Mod({\OLId{latin-looped}{L}}){!E}=\Mod({\OLId{latin-looped}{L}}){!D}$، وبيان الوجهين كما يأتي. إن كان
$\Struct{N} \models_{\OLId{latin-looped}{L}} !D$، وجد $\Struct{M} \models_{\OLId{latin-looped}{L}} !E$ بحيث
$\Struct{N} \models !D_{\Struct{M}}$، ومن ثم
$\Struct{N} \models_{\OLId{latin-looped}{L}} !E$ أيضًا (بفرض اللمة). وأما العكس، فإن كان
$\Struct{N} \models_{\OLId{latin-looped}{L}} !E$، كان $!D_\Struct{N}$ أحد حدود الانفصال في
$!D$، ولما كان $\Struct{N} \models !D_\Struct{N}$، كان كذلك
$\Struct{N} \models_{\OLId{latin-looped}{L}} !D$.
\end{proof}

\begin{thm}[مبرهنة ليندستروم]
  \ollabel{thm:lindstrom} إذا كان $\tuple{{\OLId{latin-looped}{L}},\models_{\OLId{latin-looped}{L}}}$ منطقًا مجردًا سويًا
  يجمع خاصيتي التراص ولوفنهايم--سكولم، فإن
  $\tuple{{\OLId{latin-looped}{L}},\models_{\OLId{latin-looped}{L}}} \le \tuple{{\OLId{latin-looped}{F}},\models}$ (ومن ثم يكون
  $\tuple{{\OLId{latin-looped}{L}},\models_{\OLId{latin-looped}{L}}}$ مكافئًا لمنطق الرتبة الأولى).
\end{thm}

\begin{proof}
ترجع المسألة، بمقتضى \olref{lem:lindstrom}، إلى إثبات أنه لكل
$!E \in {\OLId{latin-looped}{L}}(\Lang{L})$، حيث $\Lang{L}$ منتهية، يوجد ${\OLId{latin-ordinary}{n}} \in \Nat$ بحيث
إنه لأي !!{structure}s اثنتين $\Struct{M}$ و$\Struct{N}$: إذا كان
$\Struct{M} \equiv_{\OLId{latin-ordinary}{n}} \Struct{N}$، اتفقت $\Struct{M}$ و$\Struct{N}$
على~$!E$. فعندئذ تكافئ $!E$~!!{sentence} من الرتبة الأولى، فيترتب على ذلك
$\tuple{{\OLId{latin-looped}{L}},\models_{\OLId{latin-looped}{L}}} \le \tuple{{\OLId{latin-looped}{F}},\models}$. وحيث إن البحث في
!!{language} منتهية علائقية محضة، تجيز لنا
\olref[bas][pis]{thm:b-n-f} إحلال القضية الجبرية المقابلة
${\OLId{latin-looped}{I}}_{\OLId{latin-ordinary}{n}}(\emptyset,\emptyset)$ محل القضية
$\Struct{M} \equiv_{\OLId{latin-ordinary}{n}} \Struct{N}$.

فلنأخذ $!E$، ولنفرض خلاف المطلوب: أنه لكل ${\OLId{latin-ordinary}{n}}$ توجد
!!{structure}s~$\Struct{M}_{\OLId{latin-ordinary}{n}}$ و$\Struct{N}_{\OLId{latin-ordinary}{n}}$ بحيث
${\OLId{latin-looped}{I}}_{\OLId{latin-ordinary}{n}}(\emptyset,\emptyset)$، لكن، مثلًا،
$\Struct{M}_{\OLId{latin-ordinary}{n}} \models_{\OLId{latin-looped}{L}} !E$ و$\Struct{N}_{\OLId{latin-ordinary}{n}} \not\models_{\OLId{latin-looped}{L}} !E$. وتجيز لنا خاصية
التشاكل أن نفرض أن جميع الـ$\Struct{M}_{\OLId{latin-ordinary}{n}}$ تفسر ثوابت اللغة
بالأشياء نفسها؛ ثم إن اللغة لا تحتوي إلا عددًا منتهيًا من
الـ!!{sentence}s الذرية؛ فلنا أن نفترض اتفاق هذه البنى في الـ!!{sentence}s الذرية
التي تحققها، وإلا اقتصرنا على متتالية جزئية من الـ$\Struct{M}$. ولتكن
$\Struct{M_\infty}$ اتحاد جميع الـ$\Struct{M}_{\OLId{latin-ordinary}{n}}$، أي الـ!!{structure}
الصغرى الوحيدة التي تحتوي كل $\Struct{M}_{\OLId{latin-ordinary}{n}}$ بنية جزئية. وكما في برهان
\olref[lsp]{thm:abstract-p-isom}، لتكن $\Struct{M_\infty}^*$ توسيع
$\Struct{M_\infty}$ ذي الـ!!{domain}
$\Domain{{\OLId{latin-looped}{M}}_\infty} \cup \Domain{{\OLId{latin-looped}{M}}_\infty}^{<\omega}$، في
الـ!!{language} الموسعة التي تضم محمولي الوصل ${\OLId{latin-looped}{P}}$ و${\OLId{latin-looped}{Q}}$.

ونجري البناء نفسه لتعريف $\Struct{N}_{\OLId{latin-ordinary}{n}}$ و$\Struct{N_\infty}$ و
$\Struct{N_\infty}^*$. ولتكن الآن $\Struct{U}$ هي الـ!!{structure} التي
يضم !!{domain}ها !!{domain}s~$\Struct{M_\infty}^*$ و
$\Struct{N_\infty}^*$، وكذلك الأعداد الطبيعية~$\Nat$ مع ترتيبها
الطبيعي~$\le$، في الـ!!{language} ذات محمولات إضافية تمثل الـ!!{domain}s
$\Domain{{\OLId{latin-looped}{M}}_\infty}$ و$\Domain{{\OLId{latin-looped}{N}}_\infty}$ و
$\Domain{{\OLId{latin-looped}{M}}_\infty}^{<\omega}$ و$\Domain{{\OLId{latin-looped}{N}}_\infty}^{<\omega}$، فضلًا
عن محمولات ترمّز مجالات $\Struct{M}_{\OLId{latin-ordinary}{n}}$ و$\Struct{N}_{\OLId{latin-ordinary}{n}}$ بمعنى أن:
\begin{align*}
  \Domain{{\OLId{latin-looped}{M}}_{\OLId{latin-ordinary}{n}}} & = \Setabs{{\OLId{latin-ordinary}{a}} \in \Domain{{\OLId{latin-looped}{M}}_\infty}}{{\OLId{latin-looped}{R}}({\OLId{latin-ordinary}{a}}, {\OLId{latin-ordinary}{n}})}; &
  \Domain{{\OLId{latin-looped}{N}}_{\OLId{latin-ordinary}{n}}} & = \Setabs{{\OLId{latin-ordinary}{a}} \in \Domain{{\OLId{latin-looped}{N}}_\infty}}{{\OLId{latin-looped}{S}}({\OLId{latin-ordinary}{a}},{\OLId{latin-ordinary}{n}})}; \\
  \Domain{{\OLId{latin-looped}{M}}_\infty}^{<\omega}_{\OLId{latin-ordinary}{n}} & = \Setabs{{\OLId{latin-ordinary}{a}} \in \Domain{{\OLId{latin-looped}{M}}_\infty}^{<\omega}}{{\OLId{latin-looped}{R}}({\OLId{latin-ordinary}{a}},{\OLId{latin-ordinary}{n}})}; &
  \Domain{{\OLId{latin-looped}{N}}_\infty}^{<\omega}_{\OLId{latin-ordinary}{n}} & = \Setabs{{\OLId{latin-ordinary}{a}} \in \Domain{{\OLId{latin-looped}{N}}_\infty}^{<\omega}}{{\OLId{latin-looped}{S}}({\OLId{latin-ordinary}{a}},{\OLId{latin-ordinary}{n}})}.
\end{align*}
وللـ!!{structure}~$\Struct{U}$ أيضًا علاقة ثلاثية~${\OLId{latin-looped}{J}}$ بحيث
${\OLId{latin-looped}{J}}({\OLId{latin-ordinary}{n}},\mathbf{a},\mathbf{b})$ إذا وفقط إذا كان
${\OLId{latin-looped}{I}}_{\OLId{latin-ordinary}{n}}(\mathbf{a},\mathbf{b})$.

ويمكن الآن صوغ !!a{sentence}~$!D$ من المنطق~${\OLId{latin-looped}{L}}$، في الـ!!{language}
$\Lang{L}$ الموسعة بـ${\OLId{latin-looped}{R}}$ و${\OLId{latin-looped}{S}}$ و${\OLId{latin-looped}{J}}$ وغيرها، تقول إن $\le$ ترتيب خطي
منفصل له عنصر أول ولا عنصر أخير، وإن
$\Struct{M}_{\OLId{latin-ordinary}{n}} \models_{\OLId{latin-looped}{L}} !E$ و$\Struct{N}_{\OLId{latin-ordinary}{n}} \not\models_{\OLId{latin-looped}{L}} !E$، وإنه لكل
${\OLId{latin-ordinary}{n}}$ في الترتيب يكون ${\OLId{latin-looped}{J}}({\OLId{latin-ordinary}{n}},\mathbf{a},\mathbf{b})$ إذا وفقط إذا كان
${\OLId{latin-looped}{I}}_{\OLId{latin-ordinary}{n}}(\mathbf{a},\mathbf{b})$.

نلحق باللغة ثابتًا جديدًا~${\OLId{latin-ordinary}{n}}^*$، ثم نضم إلى $!D$ نمطًا يقتضي أن
$\bar {\OLId{latin-ordinary}{n}} < {\OLId{latin-ordinary}{n}}^*$ لكل ${\OLId{latin-ordinary}{n}} \in \Nat$ قياسي (وبصورة مكافئة، يفرض لكل عدد
قياسي~${\OLId{latin-ordinary}{n}}$ أن تحت ${\OLId{latin-ordinary}{n}}^*$ ما لا يقل عن ${\OLId{latin-ordinary}{n}}$ عنصرًا). ولا يعسر تحقيق أي جزء منتهٍ من هذا
النمط مع $!D$ في $\Struct{U}$؛ إذ يكفي اختيار عنصر قياسي بالغ الكبر.
فتقتضي خاصية التراص وجود نموذج~$\Struct{V}$ لـ$!D$ وللنمط بأسره، وفي
ترتيبه عنصر~${\OLId{latin-ordinary}{n}}^*$ يعلو كل عنصر قياسي. وبوجه خاص، تحتوي
$\Struct{V}$~!!{structure}s جزئيتين $\Struct{M_{n^*}}$ و
$\Struct{N_{n^*}}$ بحيث $\Struct{M_{n^*}} \models_{\OLId{latin-looped}{L}} !E$ و
$\Struct{N_{n^*}} \not\models_{\OLId{latin-looped}{L}} !E$. وحينئذ نعرّف مجموعة
$\mathcal{I}$ من أزواج متتاليات طولها~${\OLId{latin-ordinary}{k}}$ من $\Domain{{\OLId{latin-looped}{M}}_{{\OLId{latin-ordinary}{n}}^*}}$ و
$\Domain{{\OLId{latin-looped}{N}}_{{\OLId{latin-ordinary}{n}}^*}}$، بوضع
$\tuple{\mathbf{a},\mathbf{b}} \in \mathcal{I}$ إذا وفقط إذا كان
${\OLId{latin-looped}{J}}({\OLId{latin-ordinary}{n}}^*-{\OLId{latin-ordinary}{k}},\mathbf{a},\mathbf{b})$، حيث ${\OLId{latin-ordinary}{k}}$ طول $\mathbf{a}$ و
$\mathbf{b}$. وحيث إن ${\OLId{latin-ordinary}{n}}^*$ يعلو كل عدد قياسي، لزم أن
${\OLId{latin-ordinary}{n}}^*-{\OLId{latin-ordinary}{k}}>{\OLMathNumeral{0}}$ لكل ${\OLId{latin-ordinary}{k}}$ قياسي، وتشهد المجموعة $\mathcal{I}$ أن
$\Struct{M_{n^*}} \simeq_{\OLId{latin-ordinary}{p}} \Struct{N_{n^*}}$. غير أن
\olref[lsp]{thm:abstract-p-isom} تقضي بأن $\Struct{M_{n^*}}$ مكافئة لـ
$\Struct{N_{n^*}}$ في~${\OLId{latin-looped}{L}}$؛ فتناقض اختلافهما المفروض، وثبت المطلوب.
\end{proof}

\end{document}
